\documentclass[11pt,a4paper]{article}
\usepackage[margin=2.4cm]{geometry}
\usepackage{amsmath,amssymb,amsthm,mathtools}
\usepackage{microtype}
\usepackage{enumitem}
\usepackage[colorlinks=true,linkcolor=blue!50!black]{hyperref}
\usepackage{parskip}
\setlength{\parskip}{0.55em}
\setlength{\parindent}{0pt}

\newcommand{\R}{\mathbb{R}}
\newcommand{\X}{\mathcal{X}}
\newcommand{\Erho}{E_\rho}
\newcommand{\Frho}{F_\rho}
\newcommand{\rhad}{\widehat{\rho}}
\newcommand{\rhodelta}{\rho_\delta}
\newcommand{\Bstar}{B^{*}}
\newcommand{\Fr}{\mathcal{A}}
\newcommand{\deltaT}{\delta_T}
\newcommand{\deltaFK}{\delta_{FK}}
\newcommand{\dX}{d_{\X}}

\title{\textbf{Plan 2, the core}\\[3pt]
\large What each piece controls, and how they assemble}
\author{}\date{}

\begin{document}\maketitle\vspace{-2em}

\section*{Setup}

We want to prove $\Fr(\Omega)^2 \le C\,\deltaT(\Omega)$ (the Saint--Venant version of stability). Kohler--Jobin and bounded reduction will then upgrade this to the Faber--Krahn inequality with the right rate.

Tools we use:
\begin{itemize}[leftmargin=*,itemsep=1pt]
\item the \textbf{torsion function} $u_\Omega$: $-\Delta u = 1$ in $\Omega$, $u = 0$ on $\partial\Omega$;
\item the \textbf{superlevel sets} $\Erho = \{u_\Omega > t(\rho)\}$, parametrised so that $|\Erho| = \omega_n \rho^n$, $\rho \in [0,1]$;
\item the \textbf{rescaled sets} $\Frho = \rho^{-1}\Erho$, all of volume $|B_1| = \omega_n$;
\item the metric space $(\X, \dX)$ of finite-perimeter sets modulo translation, with $\dX([A],[B]) = \inf_v |A \triangle (B+v)|$.
\end{itemize}

As $\rho$ runs from $\rho_*$ (small) up to $1$, the family $[\Frho]$ is a path in $\X$. Our target ball $B_1$ also lives in $\X$. We will show that this path passes within distance $\sqrt{\deltaT}$ of $[B_1]$ at a specific $\rhad$ near $1$, and then transfer the closeness back to $\Omega$.

\section*{1.\ The deficit identity: two defects per level}

The starting fact is an identity (not just an inequality) that breaks the global deficit into a $\rho$-integral of two non-negative pointwise quantities:
\[
\boxed{\deltaT(\Omega) \;=\; \int_{\rho_*}^{1} \bigl(D_H(\rho) + D_I(\rho)\bigr)\, w(\rho)\, d\rho,}
\]
with explicit weight $w(\rho) > 0$ and
\[
D_H(\rho) \;:=\; \alpha(\rho)\beta(\rho) - P(\Erho)^2,\qquad
D_I(\rho) \;:=\; P(\Erho)^2 - c_n |\Erho|^{2-2/n},
\]
where $\alpha = \int_{\partial^*\Erho}|\nabla u|^{-1}$ and $\beta = \int_{\partial^*\Erho}|\nabla u|$. Both are $\ge 0$ (Cauchy--Schwarz / isoperimetric). The identity comes from refusing to discard either defect inside Talenti's symmetrisation proof of $T(\Omega) \le T(\Bstar)$.

\textbf{Consequence.} Both
\[
\int D_H\, w\, d\rho \;\le\; \deltaT, \qquad \int D_I\, w\, d\rho \;\le\; \deltaT.
\]

What each defect \emph{is}, at one level:
\begin{itemize}[leftmargin=*,itemsep=1pt]
\item $D_I(\rho)$ measures how \textbf{non-ball} the level set $\Erho$ itself is. $D_I = 0$ iff $\Erho$ is a ball.
\item $D_H(\rho)$ measures how \textbf{non-uniformly $|\nabla u|$} distributes on $\partial^* \Erho$. $D_H = 0$ iff $|\nabla u|$ is constant on the level set.
\end{itemize}

These are the only two ingredients we extract from $\deltaT$. Plan~2 uses each for one specific job. \textbf{We cannot do either job without the other.} The next two sections explain what each controls and the section after that explains why both are needed simultaneously.

\section*{2.\ $D_I$ controls how far each level is from a ball (\emph{potential})}

For each $\rho$ where $D_I(\rho)$ is below a fixed threshold (these are the \textbf{good levels}; the bad ones have small total $\mu$-measure by Markov, so we ignore them), Fusco--Julin's quantitative isoperimetric inequality says: there is a centre $z_\rho \in \R^n$ such that
\[
|\Erho \triangle B_\rho(z_\rho)|^2 \;\le\; \rho^{2n}\cdot C\, D_I(\rho).
\]
Rescaling by $\rho^{-1}$:
\[
\boxed{\dX([\Frho],[B_1])^2 \;\le\; C\, D_I(\rho).}
\]

In words: \textbf{$D_I(\rho)$ is the pointwise distance$^2$ from $[\Frho]$ to $[B_1]$ in $\X$}, up to constants.

Think of $\dX([\Frho],[B_1])^2$ as a \textbf{potential energy} for the path $[\Frho]$ in $\X$: it tells us, at each $\rho$, how far the path currently is from the target $B_1$. Integrating against $w\, d\rho$:
\[
\int \dX([\Frho],[B_1])^2 \, w\, d\rho \;\le\; C \int D_I \, w\, d\rho \;\le\; C\,\deltaT.
\]

\textbf{Could we just use this?} No. ``$\int (\text{potential}) \le \deltaT$'' means the path is on average near $B_1$ --- but we need a \emph{specific} $\rho$ where the path is provably near $B_1$, and that specific $\rho$ has to be close to $1$ (so the level set is close to $\Omega$). A pure measure statement does not let us locate a specific good $\rho$. For that we need to know the path can't dart in and out of a neighbourhood of $B_1$ --- i.e., a control on its \emph{speed}.

\section*{3.\ $D_H$ controls how fast each level moves (\emph{kinetic})}

The path $\rho \mapsto [\Frho]$ in $\X$ has a well-defined metric derivative $|\dot \Frho|_\X$ (absolutely continuous in $\rho$). Its size is determined by the normal velocity of the boundary $\partial^* \Erho$ under $\rho$. The normal velocity is
\[
V_\rho(x) \;=\; \frac{-t'(\rho)}{|\nabla u(x)|}, \qquad x \in \partial^* \Erho.
\]
After accounting for the rescaling $\Frho = \rho^{-1}\Erho$, the speed in $\X$ comes out to
\[
\boxed{|\dot \Frho|_\X^2 \;\le\; C\, D_H(\rho) \cdot (\text{explicit weight}).}
\]

Why $D_H$? Because $|\dot \Frho|_\X$ is sensitive only to the \emph{non-uniform} part of the motion of the boundary: a level set that moves uniformly outward (constant $V_\rho$) just rescales, and is invisible in $\X$ (where we mod out by translation and the volume is fixed by the parametrisation). The non-uniformity of $V_\rho$ is exactly the non-uniformity of $|\nabla u|$ on $\partial^*\Erho$, which is what $D_H$ measures.

Integrating:
\[
\int |\dot \Frho|_\X^2 \, w\, d\rho \;\le\; C \int D_H\, w\, d\rho \;\le\; C\,\deltaT.
\]

In words: \textbf{the kinetic energy of the path $[\Frho]$ in $\X$ is bounded by $\deltaT$.} If $\deltaT$ is small, the path moves slowly.

\section*{4.\ Why we need both: the abstract metric trace}

Now combine. From Sections~2 and 3:
\[
\boxed{
\int_{\rho_*}^{1}
\Bigl(
\underbrace{\dX([\Frho],[B_1])^2}_{\text{potential}}
\;+\;
\underbrace{|\dot \Frho|_\X^2}_{\text{kinetic}}
\Bigr)\, w\, d\rho
\;\le\;
C\,\deltaT(\Omega).
}
\]
This is the integrated ``action'' of the path $[\Frho]$.

\textbf{The abstract weighted metric trace lemma.} In any metric space, if an absolutely continuous path $\gamma: [a,b] \to \X$ satisfies
\[
\int_a^b \bigl(d(\gamma(\rho), p)^2 + |\dot\gamma(\rho)|^2\bigr)\, w(\rho)\, d\rho \;\le\; \varepsilon,
\]
then on any sub-window $[c,c+L]$ with weighted length $L$ there exists a specific point $\rhad \in [c, c+L]$ with
\[
d(\gamma(\rhad), p)^2 \;\le\; \frac{C \varepsilon}{L}.
\]
The proof is averaging plus the triangle inequality with the kinetic bound: \emph{Markov} on the potential gives some good level inside the window; the \emph{kinetic} term lets us transport that good level to any chosen target inside the window without spoiling the bound.

\textbf{We need both bounds.} Without the kinetic bound (only the potential bound), Markov alone gives ``most levels are near $B_1$'' but does not let us point to a specific $\rhad$ close to a chosen $\rho$. Without the potential bound (only the kinetic bound), we know the path is slow but have no anchor saying it ever comes close to $B_1$.

\textbf{Choosing the window.} We pick $L \asymp \sqrt{\deltaT}$ around the target value $\rhodelta := 1 - \kappa\sqrt{\deltaT}$. This window is short enough that $\rhad$ lands within $O(\sqrt{\deltaT})$ of $1$ --- needed for the transfer step --- and long enough that the abstract lemma still delivers $d^2 \le C \cdot \deltaT$:
\[
\boxed{
\exists\, \rhad \in [\rhodelta - c_0\sqrt{\deltaT},\, \rhodelta] \quad \text{with} \quad
\dX([F_{\rhad}], [B_1])^2 \;\le\; C\,\deltaT(\Omega).
}
\]

\textbf{Where does the radial trace lemma fit?} The FJ inequality directly delivers symmetric-difference control, but the kinetic estimate in Section 3 is written in terms of a radial quantity on $\partial^* \Erho$ (how far each boundary point is from the sphere of radius $\rho$ around $z_\rho$). The radial trace lemma is the geometric identity that converts the two FJ-controlled quantities (symmetric difference + normal oscillation, both $D_I$-controlled) into the radial form that the kinetic estimate consumes. It's a bridge, not a standalone idea: it makes Sections 2 and 3 compatible, so they can be added.

\section*{5.\ Transfer to $\Omega$, square at the end}

We now have $\dX([F_{\rhad}], [B_1])^2 \le C \deltaT$ and $\rhad \in [1 - O(\sqrt{\deltaT}), 1]$. Two consequences:

(a) Unrescale: $\Fr(E_{\rhad}) \le \dX([F_{\rhad}], [B_1]) \le \sqrt{C\deltaT}$.

(b) Boundary layer is thin: $|\Omega \setminus E_{\rhad}| = \omega_n(1 - \rhad^n) \le C\sqrt{\deltaT}$.

Triangle inequality on Fraenkel asymmetry:
\[
\Fr(\Omega) \;\le\; \Fr(E_{\rhad}) + \frac{2|\Omega \setminus E_{\rhad}|}{|\Omega|}
\;\le\; C\sqrt{\deltaT(\Omega)}.
\]

\textbf{Squaring at the very end} preserves the sharp exponent:
\[
\boxed{\Fr(\Omega)^2 \;\le\; C(n, R)\, \deltaT(\Omega).}
\]

\textbf{Closing.} Bounded reduction (BDV Lemma 5.3) removes the $R$ dependence, giving $\Fr(\Omega)^2 \le C(n)\,\deltaT(\Omega)$ for all $\Omega$. Kohler--Jobin's inequality converts this to
\[
\boxed{\deltaFK(\Omega) \;\ge\; c_{FK}(n)\,\Fr(\Omega)^2 \qquad \text{for every open $\Omega \subset \R^n$ of finite measure.}}
\]

\section*{6.\ One-page recap}

\begin{center}
\renewcommand{\arraystretch}{1.15}
\begin{tabular}{rcl}
\textbf{Identity} & : & $\deltaT(\Omega) = \int (D_H + D_I)\, w\, d\rho$.\\[2pt]
\textbf{What is $D_I$?} & : & non-ballness of $\Erho$.\\
\textbf{What is $D_H$?} & : & non-uniformity of $|\nabla u|$ on $\partial^*\Erho$.\\[6pt]
\textbf{$D_I$ controls} & : & potential $\dX([\Frho], [B_1])^2$ (via Fusco--Julin).\\
\textbf{$D_H$ controls} & : & kinetic $|\dot \Frho|_\X^2$ (via metric first variation).\\[6pt]
\textbf{Bridge} & : & radial trace lemma makes potential and kinetic compatible.\\[6pt]
\textbf{Sum} & : & $\int (d^2 + |\dot F|^2)\, w\, d\rho \le C\,\deltaT$.\\[2pt]
\textbf{Metric trace} & : & $\exists\, \rhad \in [\rhodelta - c_0\sqrt{\deltaT},\, \rhodelta]$ with $d^2 \le C\,\deltaT$.\\[2pt]
\textbf{Transfer} & : & $\Fr(\Omega) \le \Fr(E_{\rhad}) + |\Omega \setminus E_{\rhad}|/|\Omega| \le C\sqrt{\deltaT}$.\\[2pt]
\textbf{Square} & : & $\Fr(\Omega)^2 \le C\,\deltaT(\Omega)$.\\[2pt]
\textbf{KJ} & : & $\deltaFK(\Omega) \ge c_{FK}(n)\,\Fr(\Omega)^2$.\\
\end{tabular}
\end{center}

\bigskip

\textbf{The proof in one sentence.} The torsion deficit $\deltaT$ decomposes exactly into a kinetic+potential integral along the level-set foliation, both terms small force the foliation, viewed as a path in the shape space $\X$, to pass close to the unit ball at a specific level $\rhad \approx 1 - O(\sqrt{\deltaT})$; transferring back to $\Omega$ and squaring gives the sharp $\Fr(\Omega)^2 \lesssim \deltaT$.

\end{document}
